% Part: turing-machines
% Chapter: machines-computations
% Section: halting-states

\documentclass[../../../include/open-logic-section]{subfiles}

\begin{document}

\olfileid{tur}{mac}{hal}
\olsection{ఆగే స్థితులు}

\begin{explain}
అమలు చేయడానికి నిర్దేశమేదీ లేనప్పుడే మన యంత్రాలు ఆగుతాయని నిర్వచించినప్పటికీ,
ట్యూరింగ్ యంత్రాల సాధారణ ప్రాతినిధ్యాల్లో $h \in Q$ అయ్యే ప్రత్యేక
\emph{ఆగే స్థితి}~$h$ ఉంటుంది.

ఆగే స్థితి వెనుక ఉన్న భావన సరళమైనది: యంత్రం తన పనిని పూర్తిచేసినప్పుడు
(ఇన్‌పుట్‌ను అంగీకరించడానికి సిద్ధమైనప్పుడు లేదా అవుట్‌పుట్ వ్రాయడం
ముగించినప్పుడు), అది ఆగే స్థితి~$h$లోకి వెళ్లి అక్కడ ఆగుతుంది. కొన్ని
యంత్రాల్లో రెండు ఆగే స్థితులు ఉంటాయి: ఒకటి ఇన్‌పుట్‌ను అంగీకరిస్తుంది, మరొకటి
దాన్ని తిరస్కరిస్తుంది.
\end{explain}

\begin{ex}\emph{ఆగే స్థితులు}.
ఈ భావనను స్పష్టం చేయడానికి, సరి యంత్రంలో ఒక మార్పుతో మొదలుపెడదాం. ఇన్‌పుట్
సరిసంఖ్య అయినప్పుడు యంత్రం స్థితి~$q_0$లో ఆగడానికి బదులుగా, యంత్రాన్ని ఒక
ఆగే స్థితిలోకి పంపే నిర్దేశాన్ని చేర్చవచ్చు.
\[
\begin{tikzpicture}[->,>=stealth',shorten >=1pt,auto,node distance=2.8cm,
                    semithick]
  \tikzstyle{every state}=[fill=none,draw=black,text=black]

  \node[initial, state]         (A)                     {$q_0$};
  \node[state]         (B) [right of=A] {$q_1$};
  \node[state]         (C) [below of=A] {$h$};

  \path (A) edge [bend left] node {\TMtrans{\TMstroke}{\TMstroke}{R}} (B)
             edge node {\TMtrans{\TMblank}{\TMblank}{N}} (C)
        (B) edge [loop above] node {\TMtrans{\TMblank}{\TMblank}{R}} (B)
            edge [bend left] node {\TMtrans{\TMstroke}{\TMstroke}{R}} (A);
\end{tikzpicture}
\]

ఈ ఉదాహరణను మరింత విస్తరిద్దాం. ఇన్‌పుట్ బేసి అని యంత్రం నిర్ణయించినప్పుడు
అది ఎన్నటికీ ఆగదు. చక్రీయ నిర్దేశాన్ని తిరస్కరణ స్థితి~$r$లోకి వెళ్లే
నిర్దేశంతో మార్చి, యంత్రంలో ఒక \emph{తిరస్కరణ} స్థితిని చేర్చవచ్చు.
\[
\begin{tikzpicture}[->,>=stealth',shorten >=1pt,auto,node distance=2.8cm,
                    semithick]
  \tikzstyle{every state}=[fill=none,draw=black,text=black]

  \node[initial,state]         (A)                     {$q_0$};
  \node[state]         (B) [right of=A] {$q_1$};
  \node[state]         (C) [below of=A] {$h$};
  \node[state]         (D) [below of=B] {$r$};

  \path (A) edge [bend left] node {\TMtrans{\TMstroke}{\TMstroke}{R}} (B)
             edge node {\TMtrans{\TMblank}{\TMblank}{N}} (C)
        (B) edge node {\TMtrans{\TMblank}{\TMblank}{N}} (D)
            edge [bend left] node {\TMtrans{\TMstroke}{\TMstroke}{R}} (A);
\end{tikzpicture}
\]
\end{ex}

\begin{explain}
యంత్రం కొన్ని ఇన్‌పుట్‌లను ఎప్పుడు అంగీకరిస్తుందో లేదా తిరస్కరిస్తుందో స్పష్టంగా
చూపే ఇలాంటి సందర్భాల్లో, ప్రత్యేక ఆగే స్థితిని చేర్చడం ప్రయోజనకరం. అయితే,
ప్రత్యేక ఆగే స్థితిని చేర్చడం వల్ల అదనపు గణన సామర్థ్యమేదీ రాదని గమనించడం
ముఖ్యం. అదే విధంగా, ఆగడానికి తక్కువ ఆచారబద్ధమైన భావనకు దాని సొంత ప్రయోజనాలు
ఉన్నాయి. ఈ అధ్యాయంలో ఇంతవరకు ఉపయోగించిన ఆగడం యొక్క నిర్వచనం
\emph{ఆగే సమస్య} నిరూపణను అంతర్బోధాత్మకంగానూ సులభంగా ప్రదర్శించగలిగేలా
చేస్తుంది. ఈ కారణంగా మన మొదటి నిర్వచనాన్నే కొనసాగిస్తాం.
\end{explain}

\end{document}
